\uuid{cAvt}
\exo7id{6790}
\titre{exo7 6790}
\auteur{gijs}
\organisation{exo7}
\datecreate{2011-10-16}
\isIndication{false}
\isCorrection{false}
\chapitre{Théorème du point fixe}
\sousChapitre{Théorème du point fixe}
\cpge{mpsi-mp2i.fonctions-valeurs-reelles-complexes.b, pcsi.techniques-calcul-fonctions.b, ptsi.techniques-calcul-fonctions.b, mpsi-mp2i.nombres-reels-suites.d, pcsi.nombres-reels-suites.c, ptsi.nombres-reels-suites.c, mpsi-mp2i.limites-continuite.c, pcsi.limites-continuite.c, ptsi.limites-continuite.c, bcpst1.suites-reelles, bcpst1.limites-continuite.c, bcpst1.derivation.a, mp-mpi.topologie-espaces-normes.b, pc.espaces-vectoriels-normes.b, psi.espaces-vectoriels-normes.b}
\module{Topologie}
\niveau{L3}
\difficulte{}

\contenu{
\texte{
Soit $I\subset \Rr$ un intervalle fermé borné et
soit $f: I \to I$ une application dérivable verifiant
$\forall x\in I\ :\ |f'(x)| < 1$. Pour un $x_0 \in I$ on
définit la suite récurrente $a : \Nn \to I$ par
$a_0 = x_0$, $a_{n+1} = f(a_n)$. On vous demande de
montrer qu'il existe un unique $\ell \in I$, indépendant
de $x_0$, tel que $ \lim_{n\to \infty} a_n = \ell$.
Les questions suivantes peuvent vous guider dans la
démonstration.
}
\begin{enumerate}
    \item \question{Montrer que $f$ admet un point fixe unique $\ell$.}
    \item \question{Montrer que $d(a_n,\ell)$ converge.}
    \item \question{Montrer que la suite $a$ admet une sous-suite
convergente $b$.}
    \item \question{Notons $ \lim_{k\to \infty} b_k = \beta$
et  $ \lim_{n\to \infty} d(a_n,\ell) = r$.
Montrer~: $d(\beta,\ell) = r = d(f(\beta),\ell)$.}
\end{enumerate}
}
